\append{Phase spectrum of the potential $U(\mathbf{r})$}
\label{app:potential-U-min-phase}

Let the phase spectrum \cite[e.g.,][ p. 261]{blakely1996} of  
$\mathcal{F} [ U(\mathbf{r}^{\dagger}) ]$ (equation \ref{eq:potential-U-Fourier}) 
on a horizontal plane $z^{\dagger}_{0}$ above the sources be defined as follows:
\begin{equation}
	\Theta [ U(\mathbf{r}^{\dagger}) ] = \tan^{-1} \left( 
	\frac{\text{Im} \, \mathcal{F} [ U(\mathbf{r}^{\dagger}) ]}{\text{Re} \, \mathcal{F} [ U(\mathbf{r}^{\dagger}) ]} 
	\right) \: ,
	\label{eq:potential-U-phase-spectrum}
\end{equation}
where ``$\text{Re}$'' and ``$\text{Im}$'' denote the real and imaginary parts, respectively.


According to equation \ref{eq:potential-U-phase-spectrum}, the phase spectrum of $U(\mathbf{r}^{\dagger})$ (equation \ref{eq:potential-U-Fourier}) 
represents a weighted sum of the Fourier transforms $\mathcal{F} [ \sigma(\mathbf{r}') ]$
(equation \ref{eq:sigma-Fourier}) over all depth slices crossing the sources,
where the weights are real functions that exponentially decay with depth $z'$.
It means that $\Theta [ U(\mathbf{r}^{\dagger}) ]$ attributes more weight to the shallow than to the deeper parts of the sources. 
Because of that, the phase spectrum of $U(\mathbf{r}^{\dagger})$ approximates that of $\sigma(\mathbf{r}')$ and, consequently, 
$U(\mathbf{r}^{\dagger})$ has approximately the same horizontal spatial distribution of the (shallow parts of) true sources.
The same reasoning can be applied to the phase spectra of any vertical integral or vertical derivative of the potential $U$.